% Janmitra prototype-stage report. This is the only required source file.
% Compile twice with pdfLaTeX, or once with Tectonic (which performs reruns).
% The chapter structure follows the supplied TODO template. Standard report is
% used because tietreport.cls is not included in the repository. Bibliographic
% entries and the architecture figure are embedded; no external assets are needed.
\documentclass[11pt,a4paper,oneside]{report}

\usepackage[T1]{fontenc}
\usepackage[utf8]{inputenc}
\usepackage{lmodern}
\usepackage[margin=1in,headheight=15pt]{geometry}
\usepackage{microtype}
\usepackage{enumitem}
\usepackage{array,booktabs,tabularx,longtable}
\usepackage{xcolor}
\usepackage{amsmath}
\usepackage{fancyhdr}
\usepackage{xurl}
\usepackage[hidelinks]{hyperref}

\hypersetup{
  pdftitle={Janmitra: Prototype-Stage Project Report},
  pdfauthor={Dhruv Srivastava; Harkamal Singh Lubana; Paras},
  pdfsubject={UCS503P prototype implementation and evaluation}
}
\urlstyle{tt}
\setlength{\parindent}{0pt}
\setlength{\parskip}{0.45em}
\setlength{\emergencystretch}{3em}
\renewcommand{\arraystretch}{1.18}
\setlist[itemize]{leftmargin=1.5em,itemsep=0.2em,topsep=0.3em}
\setlist[enumerate]{leftmargin=1.8em,itemsep=0.25em,topsep=0.3em}
\newcolumntype{Y}{>{\raggedright\arraybackslash}X}
\newcolumntype{P}[1]{>{\raggedright\arraybackslash}p{#1}}
\newcommand{\code}[1]{\texttt{\detokenize{#1}}}
\newcommand{\ReportDate}{10 October 2026}
\newcommand{\Revision}{e34544d}
\newcommand{\Advisor}{Dr Paramveer Kaur}

\pagestyle{fancy}
\fancyhf{}
\fancyhead[L]{\small Janmitra}
\fancyhead[R]{\small Prototype-Stage Report}
\fancyfoot[C]{\thepage}
\renewcommand{\headrulewidth}{0.4pt}
\fancypagestyle{plain}{%
  \fancyhf{}\fancyfoot[C]{\thepage}\renewcommand{\headrulewidth}{0pt}}

\begin{document}

\hypersetup{pageanchor=false}
\begin{titlepage}
  \centering
  {\large\bfseries Thapar Institute of Engineering and Technology\par}
  \vspace{0.2cm}
  {\large Patiala\par}
  \vspace{0.9cm}
  {\large UCS503P --- Course Project\par}
  \vspace{1.1cm}
  {\Huge\bfseries Janmitra\par}
  \vspace{0.4cm}
  {\LARGE A Voice-First Civic Scheme Guidance\par
    Platform for Rural India\par}
  \vspace{0.8cm}
  {\Large\bfseries Prototype-Stage Project Report\par}
  \vspace{0.2cm}
  {\large Mid-Semester Evaluation\par}
  \vspace{1.2cm}
  {\large Submitted to\par}
  \vspace{0.15cm}
  {\Large\bfseries \Advisor\par}
  \vspace{1.0cm}
  {\large Submitted by\par}
  \vspace{0.3cm}
  \begin{tabular}{@{}ll@{}}
    \textbf{Student} & \textbf{Roll number} \\
    \midrule
    Dhruv Srivastava & \texttt{1024170394} \\
    Harkamal Singh Lubana & \texttt{1024170396} \\
    Paras & \texttt{1024170395}
  \end{tabular}
  \vspace{0.4cm}
  \par{\small
    \texttt{dsrivastava\_be24@thapar.edu}\par
    \texttt{hlubana\_be24@thapar.edu}\par
    \texttt{pparas\_be24@thapar.edu}\par}
  \vfill
  {\large\bfseries Computer Science and Engineering Department\par}
  \vspace{0.4cm}
  {\large \ReportDate\par}
\end{titlepage}

\pagenumbering{roman}
\hypersetup{pageanchor=true}
\chapter*{Abstract}
\addcontentsline{toc}{chapter}{Abstract}
Janmitra is a voice-first civic-guidance project intended to help rural citizens
discover government support, understand eligibility conditions, prepare documents,
and request human assistance. This report describes the implemented prototype at
repository revision \texttt{\Revision}. The current interaction channel is browser
audio over LiveKit, with Gemini Live handling native speech input and output.
A FastAPI backend provides structured scheme retrieval, deterministic eligibility
evaluation, document checklists, versioned publication, an operator handoff queue,
and conversation and audit records.

The prototype includes 18 structured scheme drafts across six categories. All
remain pending human review; five contain draft rule sets. Retrieval can supply
these references with their review status preserved, while verified eligibility
and checklist tools require published records. A local backend verification run
completed with 310 passing tests and five skipped PostgreSQL tests; two test
modules requiring optional voice-runtime dependencies were excluded. Lint checks
passed. Repository documentation also records one successful synthetic English
voice call, alongside intermittent failures and unresolved answer-quality gaps.

The work establishes a testable application foundation and a browser voice
prototype. Phone/SIP access from the original proposal, complete source ingestion,
human verification of the catalogue, representative multilingual evaluation,
measured load scaling, and staging deployment remain outstanding. This report
distinguishes implemented mechanisms, observed results, and future objectives.

\tableofcontents
\clearpage
\pagenumbering{arabic}

% ================================================================
\chapter{Introduction}

\section{Background and Context}
Government welfare information is distributed across departmental websites,
guidelines, notices, and application portals. A citizen may describe a practical
need---such as support for farming, employment, housing, insurance, or a small
business---without knowing the relevant scheme name or department. Formal
language, unfamiliar terminology, and multi-page navigation can make the
information difficult to use even when it is publicly available.

The higher-order goal of Janmitra is to improve access to this information through
conversation. The original proposal envisaged a service reachable by an ordinary
phone call, allowing citizens to speak in a familiar language or dialect without
installing an application \cite{proposal}. The prototype develops the speech and
backend workflow through a browser voice channel. This provides a practical
environment for integration and demonstration, while retaining phone accessibility
as an unmet objective of the original proposal.

\subsection{Problem Statement}
The project addresses five connected difficulties:
\begin{enumerate}
  \item \textbf{Fragmented information:} relevant facts may be spread across
    multiple official sources and versions.
  \item \textbf{Language and literacy barriers:} a citizen may be more comfortable
    describing a need aloud than reading formal scheme descriptions.
  \item \textbf{Discovery difficulty:} a person may know their circumstances but
    not the category or name of a suitable programme.
  \item \textbf{Eligibility and document uncertainty:} general descriptions do not
    always explain which conditions apply to a particular person or what they
    should prepare before visiting an office.
  \item \textbf{Incomplete assistance:} a self-service interaction needs an
    explicit route to human help when the issue cannot be resolved.
\end{enumerate}

The engineering problem is to combine a natural spoken interaction with
traceable information and predictable backend decisions. Language generation
alone cannot establish personal eligibility, official approval, or completion
of an application. These distinctions guide the prototype's tool boundaries.

\section{Project Objectives}
The prototype checkpoint is organised around the following observable objectives.
Their implementation and validation status are examined in Chapter~3.
\begin{enumerate}
  \item \textbf{Demonstrate the voice-to-tool path:} connect browser audio to
    LiveKit and Gemini Live, retrieve scheme context through the API, and return
    a spoken response during a controlled demonstration.
  \item \textbf{Define a reusable catalogue:} represent scheme descriptions,
    citations, review state, documents, application steps, and optional rules
    using one validated schema. Include an initial, extensible reference set.
  \item \textbf{Implement deterministic guidance:} evaluate rule conditions and
    conditional document requirements with explicit outcomes and boundary tests.
  \item \textbf{Preserve accountable publication:} require verified records for
    publication, retain previous versions, and associate verification with a
    named reviewer and date.
  \item \textbf{Support human assistance:} create a queued handoff through a
    validated backend operation and expose operator status updates and context.
  \item \textbf{Establish repeatable verification:} provide unit and API tests,
    linting, database migrations, container definitions, and voice-evaluation
    tools before expanding pilot and performance testing.
\end{enumerate}

\subsection{Time-to-Value and Development Approach}
The proposal prioritised an early working voice loop, followed by catalogue
expansion, ingestion, and load testing. The prototype follows an incremental
approach: establish data contracts and domain services, integrate voice through
HTTP tools, add representative source records, then strengthen retrieval,
concurrency controls, and evaluation support. Each increment provides a concrete
artefact that can be inspected or tested rather than relying on a final integrated
demonstration alone.

\subsection{Prototype Scope and Changes from the Proposal}
The report covers the current repository implementation \cite{implementation},
associated technical documentation, and the local checks described in Chapter~3. It does not assume
that an externally hosted deployment or a populated production database exists.

\begin{table}[htbp]
\centering\small
\begin{tabularx}{\textwidth}{@{}P{0.23\textwidth}YY@{}}
\toprule
\textbf{Area} & \textbf{Original proposal} & \textbf{Current prototype} \\
\midrule
Citizen access & Ordinary telephone call through a SIP bridge.
  & Browser microphone and speaker through LiveKit; no telephone integration. \\
Scheme guidance & Facts obtained from a verified, published catalogue.
  & Published facts and pending references are distinguished; qualified general
    guidance is allowed when context is incomplete. \\
Eligibility & Deterministic pre-checks over reviewed rules.
  & Engine and tests implemented; bundled rule sets still await human review. \\
Source ingestion & Import, AI extraction, human review, and versioned publish.
  & Structured JSON preparation, review attestation, and versioned publication;
    complete ingestion workflow remains future work. \\
Delivery and scale & Staging pipeline and measured replica comparisons.
  & CI checks, containers, migrations, and concurrency tests; staging CD and
    load measurements remain outstanding. \\
\bottomrule
\end{tabularx}
\caption{Relationship between the proposal and the prototype checkpoint.}
\label{tab:scope}
\end{table}

The browser channel currently requires a compatible device and data connection.
It therefore does not yet deliver the proposal's access from an ordinary phone.
Similarly, reference-based guidance must not be described as a fully
human-verified knowledge base. These scope differences are material to evaluating
the prototype fairly \cite{architecture,voice-design}.

% ================================================================
\chapter{Methodology and Design}

\section{System Architecture}

\subsection{High-Level Design}
Janmitra uses a modular backend application with a separate voice worker. The
backend contains conversation, catalogue, eligibility, handoff, and audit modules.
The voice worker calls the backend through typed HTTP endpoints and does not
duplicate eligibility rules or write directly to the database.

\begin{figure}[htbp]
\centering
\setlength{\fboxsep}{8pt}
\fbox{\parbox{0.83\textwidth}{\centering
  \textbf{Browser voice channel}\\
  Citizen microphone, speaker, and call controls}}
\par\smallskip $\updownarrow$\quad WebRTC audio\par\smallskip
\fbox{\parbox{0.83\textwidth}{\centering
  \textbf{LiveKit room and Python voice worker}\\
  Gemini Live native audio session; tool invocation; transcript and timing events}}
\par\smallskip $\updownarrow$\quad Authenticated HTTP tools\par\smallskip
\fbox{\parbox{0.83\textwidth}{\centering
  \textbf{FastAPI backend}\\
  Conversation state; catalogue retrieval; deterministic eligibility;\\
  document checklists; handoff queue; publication; audit}}
\par\smallskip $\updownarrow$\quad Structured data access\par\smallskip
\fbox{\parbox{0.83\textwidth}{\centering
  \textbf{PostgreSQL and local reference catalogue}\\
  Published versions and durable records in PostgreSQL;\\
  pending scheme JSON files available as labelled reference context}}
\caption{Implemented prototype architecture. Admin and operator API clients
access the backend separately. A phone/SIP bridge is not part of this implementation.}
\label{fig:architecture}
\end{figure}

The database preserves application records across API requests. The active
real-time audio session remains the responsibility of its voice worker. This
separation supports independent growth in API and voice-worker capacity, but a
measured scaling claim requires the load experiments described later.

\subsection{Technical Stack}
\begin{itemize}
  \item \textbf{Backend language and framework:} Python 3.12 or later, FastAPI,
    and Uvicorn for asynchronous HTTP services.
  \item \textbf{Validation and persistence:} Pydantic request and record models,
    SQLAlchemy asynchronous sessions, PostgreSQL, the asyncpg driver, and Alembic
    migrations. SQLite supports isolated local tests.
  \item \textbf{Voice integration:} LiveKit Agents and its Google plugin are
    pinned to version 1.8.5. The configured speech model is
    \code{gemini-3.1-flash-live-preview}.
  \item \textbf{Model separation:} Gemini Live handles the spoken exchange;
    a separate text adapter supports backend text tasks. Mock and failure adapters
    support controlled tests without real-provider calls.
  \item \textbf{Integration and delivery tools:} HTTPX, pytest, pytest-asyncio,
    Ruff, Docker Compose, and GitHub Actions.
\end{itemize}

These are the technologies declared by the checked revision. Dependency
configuration is evidence of the intended runtime, not evidence that every
component has been deployed or benchmarked.

\section{Implementation Details}

\subsection{Voice Conversation and Session Lifecycle}
The worker connects to a LiveKit room, creates a backend conversation, and starts
an \code{AgentSession} with Gemini Live. Native audio input and output are handled
within the same model session; the implementation does not introduce a separate
speech-to-text service followed by a separate text-to-speech service.

The current flow is citizen-initiated: the citizen speaks first, and the prompt
requests a short AI disclosure in the initial spoken reply. The agent is instructed
to use the citizen's latest language, keep spoken steps understandable, and ask one
question at a time. These are behavioural instructions whose quality must still be
checked in real conversations.

The worker configures voice activity detection with 100~ms prefix padding and
500~ms silence endpointing. The default thinking level is \code{low}, configurable
within the same live model. These configuration values are not measured response
latencies. Transcripts and timing events are written asynchronously, and cleanup
drains pending writes and closes the backend conversation through a single cleanup
task. Safe status messages communicate readiness and recovery state to the client.

\subsection{Typed Backend Tools and Request Flow}
The agent exposes four primary operations:
\begin{description}[leftmargin=1.7em,style=nextline]
  \item[\code{find_service}] Retrieves relevant scheme context for the citizen's
    stated need, optionally narrowed by category.
  \item[\code{check_eligibility}] Evaluates submitted answers against the selected
    published record's deterministic rule set.
  \item[\code{get_documents}] Produces the published record's checklist, including
    conditional requirements.
  \item[\code{request_handoff}] Requests human assistance using contact details,
    an issue summary, and reported escalation signals.
\end{description}

Each tool request identifies the conversation and requested language. The backend
validates the payload, checks the caller's role, retrieves the conversation, and
performs the appropriate domain operation. Scheme facts, eligibility results,
and handoff records are returned as structured responses. The model explains
those results in speech; it cannot directly write a database publication or
substitute its own verdict for the eligibility engine's result.

\subsection{Canonical Scheme Records and the Reference Catalogue}
A \code{ServiceRecord} contains a stable slug, localised name and description,
aliases, category, benefit and eligibility summaries, optional rule set,
documents, application steps, and citation metadata. The citation identifies the
source, publisher, review state, verification date, and reviewer when verified.

The current reference directory contains 18 JSON records. Their category
distribution is shown in Table~\ref{tab:catalogue}. All 18 have
\code{pending_review} status. Five records---APY, PM-SYM, PMJJBY, PMSBY, and
PMUY---contain draft rule sets; 15 records contain nonempty document lists.
The remaining three do not yet supply document lists. These counts describe the
repository data and do not establish that any scheme has been human-verified or
published in a deployment.

\begin{table}[htbp]
\centering\small
\begin{tabularx}{\textwidth}{@{}lrY@{}}
\toprule
\textbf{Category} & \textbf{Records} & \textbf{Record identifiers} \\
\midrule
Loan & 4 & KCC, PMMY, PMEGP, PM-Vishwakarma \\
Banking & 2 & PMJDY, DAY-NRLM \\
Grant & 4 & PM-KISAN, PMAY-G, PMGKAY, PMUY \\
Insurance & 4 & PM-JAY, PMFBY, PMJJBY, PMSBY \\
Pension & 3 & APY, PM-SYM, NSAP \\
Other & 1 & Rural-employment reference with stable slug \code{mgnrega} \\
\midrule
\textbf{Total} & \textbf{18} & All pending human review \\
\bottomrule
\end{tabularx}
\caption{Inventory of bundled scheme drafts at revision \texttt{\Revision}.
Identifiers are catalogue labels, not a certification of current programme details.}
\label{tab:catalogue}
\end{table}

The catalogue review checklist requires comparison of amounts, ages, exclusions,
documents, and application steps against official sources. Automated source-check
notes assist that work but explicitly do not constitute human approval
\cite{catalogue-review,source-check}. This report consequently does not reproduce
the drafts' financial benefits or legal eligibility conditions as verified facts.

\subsection{Retrieval and Grounding}
The recent retrieval implementation combines Unicode-aware normalisation,
multilingual phrase expansion, named-scheme matching, and a curated mapping from
expressed needs to candidate schemes. It attempts to preserve distinct needs in
a compound request rather than returning several variants of the same topic.
This is lightweight retrieval over structured records; it does not use an
embedding service or a vector database.

The voice client requests \code{response_mode="context"}. The API combines
published records with local pending references, replaces a pending record when
a published version of the same slug exists, and returns at most three selected
records in total. The \code{context_order} field preserves ordering across
published \code{matches} and \code{reference_context}. Citations, service versions
where available, and verification status accompany the context.

Gemini Live then generates the spoken answer directly from this result, avoiding
an additional text-model answer-generation round trip in the voice path. A
\code{prepared_answer} mode remains for compatible HTTP clients. The current
prompt permits qualified general orientation from trained knowledge when
retrieved context is incomplete, while forbidding invented current amounts,
deadlines, confirmed eligibility, and completed actions. This broadens the original
proposal's verified-record-only guidance policy and requires explicit quality
evaluation \cite{voice-design}.

\subsection{Deterministic Eligibility and Document Checklists}
Eligibility rules are stored as data within a versioned scheme record. Questions
declare answer types and optional limits. Named conditions contain comparisons,
and optional \code{applies_when} guards determine when a condition is relevant.
The decision combines conditions through \code{all_of}, \code{any_of}, and
\code{none_of} groups.

The engine validates answers before evaluation and returns one of three states:
\code{eligible}, \code{not_eligible}, or \code{needs_more_info}. A missing answer
remains unknown, allowing the backend to supply the next relevant question.
Condition traces identify the tests involved in the result and may include
supporting source text. Numeric validation rejects inappropriate types,
non-finite values, and values outside declared bounds.

For illustration only, a synthetic rule might require an applicant to be at
least 18 years old. An answer below that boundary fails the condition; an omitted
age requests more information. This example explains engine behaviour and is not
a claim about a real scheme.

Document checklists use the same structured answers. An unconditional document
is included directly. A document whose condition is false is omitted; when the
condition is unknown, it remains visible as conditional. The verified endpoints
require a published record, and eligibility also requires an associated rule set.
Pending draft rules are not sufficient for a verified eligibility decision.

\subsection{Versioned Publication and Review Workflow}
The \code{Service} table holds scheme identity and a pointer to its current
version. \code{ServiceVersion} holds the payload, version number, publication
state, author, notes, and timestamp. Publishing appends a new version and updates
the pointer while retaining earlier versions. History and field-level diffs are
available through admin routes.

The seed command validates local JSON records. To publish pending records, the
reviewer supplies a name using \code{--actor} and explicitly attests to review
using \code{--confirm-reviewed}. The command then marks the records verified for
publication. The reviewer must actually check the sources before giving this
attestation; the flag itself cannot prove factual correctness. PostgreSQL advisory
and row locks serialize publication for a scheme, including its first publication.

The current workflow is local record preparation, human review, and CLI/API
publication. A complete official-source importer, extraction-to-draft workflow,
and integrated review dashboard are not established by this backend. A source
snapshot model and model-adapter drafting method provide foundations, but the full
ingestion process remains future work.

\subsection{Human Handoff and Operator Queue}
The handoff service evaluates reported signals such as a citizen request for a
person, an out-of-scope request, no matching scheme, low confidence, or repeated
tool problems. A successful request stores an issue summary, optional contact
name and phone number, category, language, and conversation reference. The
conversation is marked \code{handed_off}.

Operators can list queued requests, inspect the related conversation events,
add notes, and move a request through allowed states. The principal progression
is \code{new} to \code{contacted} to \code{resolved}; direct resolution from
\code{new} is also allowed. A queued handoff is a record for subsequent operator
attention. The implementation does not transfer a live telephone call or
guarantee that a callback has occurred.

\subsection{Persistence, Audit, and Access Control}
The principal database entities are services, service versions, source snapshots,
conversations, conversation events, handoff requests, and audit events. Alembic
provides an initial schema and a subsequent migration aligning structured payload
columns with PostgreSQL JSONB. Category and identity lookups have relational
indexes; current retrieval still ranks candidate records in application code,
so JSONB storage alone does not establish indexed full-catalogue search capacity.

Conversation events retain ordered transcript and tool context. Audit events
record critical actions with request identifiers, actors, and entity references.
Conversation row locks serialize related state transitions and event allocation
in PostgreSQL. Admin, operator, and voice API keys restrict access by role.
This is a prototype access-control mechanism rather than a complete user-account
and organisational permissions system.

The prompt prohibits requesting Aadhaar numbers, full bank details, OTPs, PINs,
and passwords. Contact capture is optional and associated with handoff. However,
transcripts are persisted, and prompt instructions alone cannot guarantee that
unsolicited sensitive information is never stored. Redaction, consent, retention,
and operator-access procedures require further work before a citizen pilot.

\subsection{Development, CI, and Deployment Foundations}
The backend GitHub Actions workflow installs development, voice, and Gemini
dependencies, runs Ruff and pytest, and provisions PostgreSQL 16 for integration
tests. These tests cover migrations and concurrency behaviour beyond what SQLite
can establish. Workflow configuration documents the intended checks; the local
results reported here do not substitute for a retrieved CI run result.

Docker Compose defines PostgreSQL and the API, with health checks and a persistent
database volume. The API container applies migrations before starting Uvicorn.
The voice worker is started separately. The repository provides deployment
foundations, but no completed staging CD pipeline or measured multi-replica
deployment is claimed in this report.

% ================================================================
\chapter{Results and Evaluation}

\section{Testing Methodology}

\subsection{Unit Testing}
The tests exercise rule comparisons and decision combinations, missing answers,
numeric validation, conditional documents, catalogue ranking, schema constraints,
review attestation, and handoff trigger selection. Mock adapters make expected
outputs reproducible and avoid provider cost. Synthetic test records are clearly
separate from the pending government-scheme catalogue; their conditions are test
fixtures and must not be reused as verified public guidance.

\subsection{Integration Testing}
HTTP tests run the FastAPI application against isolated SQLite databases. They
exercise role restrictions, published-record retrieval, reference context,
eligibility and document responses, handoff creation and updates, conversation
events, and audit data. PostgreSQL-specific tests additionally cover migration
round trips, model/schema agreement, concurrent publications, event numbering,
write-once guidance timestamps, and competing terminal state transitions.

The prototype also contains voice-client, agent, and worker tests. Worker tests
use controlled runtime substitutes to exercise cleanup, notifications, and
session behaviour. These tests help validate orchestration but cannot demonstrate
real speech quality or provider reliability.

\subsection{Local Verification for This Report}
A local Python 3.12 run on \ReportDate\ checked the current revision with model
calls mocked and configuration-file loading disabled. It produced
\textbf{310 passing tests and five skipped tests}, with no failed tests in that
selected run. The five skipped tests require an explicitly configured PostgreSQL
test database.

The two modules \path{tests/test_voice_agent.py} and
\path{tests/test_voice_worker.py} were excluded because the temporary verification
environment did not contain the optional LiveKit/Google runtime dependencies.
Other voice-related checks, including prompt, client, and model-setting tests,
were included. Thus, 310 passing tests describes the executed selection, not a
claim that every repository test or a live voice deployment passed.

Ruff passed over \code{app}, \code{tests}, \code{scripts}, and \code{alembic}.
Two non-failing warnings were emitted by the test run: a cache-setting warning
caused by disabling pytest's cache plugin, and a dependency deprecation warning
for an HTTP status constant. No credentials, real-provider calls, or deployed
database contents were needed for these checks.

\subsection{Voice, Performance, and User Testing}
Two scripts support separate evaluation needs:
\begin{itemize}
  \item \path{scripts/smoke_voice.py} exercises the LiveKit voice flow and checks
    audio, retrieval, response, and conversation lifecycle behaviour.
  \item \path{scripts/evaluate_voice_answer.py} supports isolated native-audio
    answer evaluation using supplied WAV files and local retrieval. It uses real
    Gemini quota but does not exercise LiveKit or the HTTP transport path.
\end{itemize}

The evaluation fixture defines 16 scenarios covering mixed needs, urgent
situations, food, health, employment, pension, education, fraud, family distress,
low literacy, multilingual speech, and follow-up context. The rubric considers
usefulness, grounding, understandable speech, one-question behaviour, handling
of multiple needs, and privacy. Unsupported material claims and invented completed
actions are hard failures. The existence of these fixtures does not mean all
scenarios have passed \cite{evaluation-fixtures}.

The proposed pilot remains a future evaluation: recruit 8--12 participants across
languages and dialects and compare two or three journeys with equivalent portal
tasks. Completion, time, errors, and 1--5 clarity/trust feedback should be recorded.
No completed participant study is claimed at this checkpoint.

\section{Performance Metrics}

\subsection{Definition and Interpretation of Timing}
The proposal defines Time-to-Guidance as:
\[
  \mathrm{TTG} = t_{\text{grounded answer received}}
                 - t_{\text{call connected}}.
\]
The current backend stores \code{connected_at} and \code{first_guidance_at}.
The latter is set when a qualifying backend tool result is produced. This is a
useful internal proxy, but it does not establish when a citizen actually heard
the answer. In the browser prototype, conversation creation also need not coincide
exactly with media connection. A proposal-level TTG measurement therefore needs
aligned connection, grounding, and audio-delivery events.

Additional instrumentation separates:
\begin{itemize}
  \item \textbf{Worker startup:} conversation creation and session startup,
    excluding initial room connection and HTTP client creation.
  \item \textbf{Retrieval time:} backend catalogue lookup and ranking duration.
  \item \textbf{Voice turn timing:} available SDK timing values, including
    time to first token, end-to-end turn latency, and playback latency.
  \item \textbf{Speech-end to speaking-state delay:} elapsed time between the
    detected end of citizen speech and the agent entering the speaking state;
    this excludes browser playout.
\end{itemize}

These clocks describe different intervals and must not be combined or relabelled
as a single end-to-end latency result. Future measurements should distinguish
cold startup from warm turns and report sample counts, medians, and slow-tail
latencies.

\subsection{Observed and Available Results}
Table~\ref{tab:results} combines locally checked implementation evidence with
clearly attributed prior observations. The voice timings come from the repository
runbook's account of one synthetic English call; they were not newly measured for
this report \cite{runbook}.

\begin{longtable}{@{}P{0.26\textwidth}P{0.21\textwidth}P{0.45\textwidth}@{}}
\caption{Prototype results and the limits of the available evidence.}
\label{tab:results}\\
\toprule
\textbf{Measure} & \textbf{Observed value/status} & \textbf{Interpretation} \\
\midrule
\endfirsthead
\toprule
\textbf{Measure} & \textbf{Observed value/status} & \textbf{Interpretation} \\
\midrule
\endhead
\bottomrule
\endfoot
Local automated tests & 310 passed; 5 skipped & Selected local suite, with two
voice-runtime modules excluded and PostgreSQL-specific tests skipped. \\
Lint & Passed & Ruff checked backend application, tests, scripts, and migrations. \\
Reference catalogue & 18 records & Six categories; all 18 pending human review. \\
Draft rule coverage & 5 records & Rule sets exist; official factual correctness
and deployment publication are not established by this count. \\
Document-list coverage & 15 of 18 records & Nonempty lists in repository drafts;
these are not yet verified checklists. \\
Voice scenarios & 16 defined & Evaluation inputs and rubric exist; no complete
scenario pass rate has been established. \\
Agent readiness & 2,431 ms & Prior runbook observation from one successful
synthetic English call; not a distribution or reliability target. \\
Speech-end to first returned audio & 3,064 ms & Same prior call; not equivalent
to call-connect-to-grounded-answer TTG or browser playout timing. \\
Median end-to-end TTG & Not established & Requires repeated real-model sessions
with aligned connection and delivered-answer timestamps. \\
Citation correctness & Not established & Citation metadata exists; semantic
claim-to-source support still requires labelled review. \\
Handoff precision/recall & Not established & Trigger logic is tested; no completed
labelled conversational evaluation is reported. \\
Throughput and replica scaling & Not established & No measured 1-, 2-, and
4-replica comparison or before/after bottleneck experiment is available. \\
User clarity and availability & Not established & Pilot feedback and repeated
availability windows remain future work. \\
\end{longtable}

\section{Analysis}

\subsection{Achievement of Prototype Objectives}
The backend now supports a complete sequence of typed operations: conversation
creation, scheme discovery, deterministic guidance over published records,
handoff creation, operator updates, and audit inspection. The selected local test
run provides evidence for these application contracts. The reference catalogue
and ranking logic also enable useful retrieval during development before the
human-publication process is completed.

The voice worker and supporting scripts establish an integration path, and the
runbook records a successful synthetic voice exchange. However, this remains
partial evidence for the voice objective. The same documentation reports a
subsequent integrated call that failed to complete guidance, and the voice-design
notes record inconsistent behaviour in complex requests. Full phone accessibility,
verified catalogue deployment, and representative speech reliability are therefore
not completed objectives.

\subsection{Recent Development Progress}
The report incorporates the latest repository history rather than carrying
forward the earlier proposal's future-tense claims as completed work:
\begin{itemize}
  \item \textbf{Commit \texttt{595fc75}:} introduced the voice integration,
    browser harness, and deployment foundations.
  \item \textbf{Commit \texttt{992d638}:} added the Week~5 backend test suite and
    accompanying verification journal.
  \item \textbf{Commit \texttt{def3630}:} expanded the catalogue and retrieval,
    refined Gemini voice behaviour and timing, strengthened concurrency and
    numeric validation, added PostgreSQL and voice tests, and updated documentation.
  \item \textbf{Commit \texttt{\Revision}:} merged the test work and resolved the
    reported test merge conflict; this is the revision used for the report.
\end{itemize}

\subsection{Challenges and Design Responses}
\textbf{Balancing helpfulness and verification.}
The initial verified-only path requires reviewed data before it can answer useful
questions. The prototype now exposes pending references with explicit review
metadata while preserving the publication gate for verified tools. This improves
development-time usefulness but increases the need to test whether spoken answers
preserve uncertainty and material conditions.

\textbf{Understanding ordinary and mixed needs.}
Users may describe several problems without naming a scheme. Multilingual phrase
expansion, curated need concepts, and combined ranking improve retrieval for such
inputs. These mappings remain finite and require extension and evaluation for
unseen phrasing, regional languages, and state-specific support.

\textbf{Managing voice latency and lifecycle.}
The context response mode removes an intermediate text-answer generation step
from the voice path. Timing events distinguish retrieval, startup, and turn delays.
Asynchronous transcript writes and single-task cleanup improve lifecycle handling.
These are implemented mechanisms; a statistically supported latency improvement
still requires comparable before-and-after runs.

\textbf{Preserving consistency under concurrency.}
Application state is persisted and critical transitions use database locks.
PostgreSQL migration and concurrency tests now exist and are configured in CI.
The local SQLite run cannot validate those PostgreSQL behaviours, so the report
does not infer concurrency capacity from its pass count.

\subsection{Limitations at the Prototype Checkpoint}
\begin{enumerate}
  \item The browser channel does not satisfy the ordinary-phone accessibility
    objective, and there is no live telephone transfer.
  \item All bundled records still require named human review. Document and rule
    coverage are incomplete, and the full source-import workflow is unfinished.
  \item Real-model answer quality and multilingual behaviour remain uneven.
    Documentation reports wrong-language responses, missed retrieval, multiple
    questions, or omitted material conditions in some trials.
  \item Integrated live sessions have shown intermittent 1011 errors. A successful
    direct SDK call does not establish the cause of an integrated failure or prove
    that the browser voice path is reliable.
  \item End-to-end TTG, citation correctness, handoff precision/recall, pilot
    feedback, and measured scalability have not yet been established.
  \item Transcript privacy requires enforced controls beyond prompt instructions,
    and staging deployment and operational monitoring remain incomplete.
\end{enumerate}

These limitations define the next engineering stage. They do not invalidate the
tested backend mechanisms, but they constrain the claims that can be made about
citizen readiness.

% ================================================================
\chapter{Conclusion}

\section{Summary of Work}
Janmitra has progressed from a proposed phone-accessible guidance system to an
implemented browser voice prototype supported by a modular backend. The work
includes typed guidance tools, a structured reference catalogue, deterministic
eligibility and document logic, preserved publication history, operator handoff
records, audit trails, migrations, containers, CI checks, and voice-evaluation
support. The latest changes substantially expand retrieval and improve the
separation between verified records and pending reference material.

The selected local verification run passed 310 tests and lint checks. The report
also identifies the unexecuted test scope and the limited nature of prior live
voice observations. This makes the checkpoint a documented implementation result
without treating future pilot or scalability targets as measured achievements.

\section{Key Findings}
\begin{itemize}
  \item \textbf{Explicit tool boundaries make guidance testable.} Eligibility,
    publication, and queued handoff creation can be checked independently of
    speech generation and provider availability.
  \item \textbf{Data review is an operational dependency.} Schema-valid JSON and
    official-source links are useful foundations, but they do not replace a
    human check of the actual scheme conditions.
  \item \textbf{Retrieval must preserve the citizen's situation.} Combined
    ranking and need-based matching help retain multiple concerns, while
    uncertainty and record provenance must survive into the spoken response.
  \item \textbf{Different evidence answers different questions.} Unit tests
    establish application behaviour; live audio trials, semantic review, pilot
    feedback, and load testing establish separate aspects of usability and capacity.
\end{itemize}

\section{Future Work and Recommendations}
\begin{enumerate}
  \item \textbf{Complete catalogue review and publication.} Review all initial
    records against official sources, complete missing checklists, verify every
    rule branch, and publish an initial trusted set with reviewer attribution.
  \item \textbf{Complete the ingestion workflow.} Add official-source import,
    retained snapshots, AI-assisted extraction into drafts, reviewer comparison,
    and controlled publication with visible history.
  \item \textbf{Validate voice quality and reliability.} Run the full 16-scenario
    set, assess transcripts and audio semantically, exercise follow-up context
    and interruptions, and investigate intermittent integrated-call failures.
  \item \textbf{Resolve the channel scope.} Implement and validate phone/SIP
    access if the original ordinary-phone objective is retained. If browser access
    becomes the final scope, explicitly revise the accessibility claims and pilot
    assumptions rather than describing it as an equivalent phone service.
  \item \textbf{Measure the proposed metrics.} Align connection and
    response-delivery events; establish TTG distributions, citation correctness,
    handoff precision/recall, clarity ratings, and availability over test windows.
  \item \textbf{Demonstrate scalability.} Build the mocked-model load harness,
    compare 1, 2, and 4 API replicas under equivalent conditions, identify the
    first bottleneck, apply a justified fix, and report before-and-after results.
    Measure voice-worker capacity separately from HTTP capacity.
  \item \textbf{Prepare a controlled pilot deployment.} Add staging CD,
    operational checks, transcript redaction and retention controls, and a
    documented operator process before recruiting the proposed 8--12 users.
\end{enumerate}

\section{Final Remarks}
The prototype demonstrates how conversational AI can be connected to structured
public-service information and deterministic backend operations. Its main
engineering contribution at this stage is the separation of speech interaction,
retrieval, trusted decisions, and persistent records into components that can be
examined and tested. Completing human data review and evaluating real citizen
journeys are the next steps toward the proposal's wider goal of accessible,
dependable civic guidance.

% All references are embedded so this deliverable remains a single .tex file.
\clearpage
\phantomsection
\addcontentsline{toc}{chapter}{Bibliography}
\begin{thebibliography}{9}

\bibitem{proposal}
Dhruv Srivastava, Harkamal Singh Lubana, and Paras.
\emph{Janmitra: A Voice-First Civic Scheme Guidance Platform for Rural India}.
UCS503P project proposal. Repository source: \path{project-proposal/main.tex}.

\bibitem{implementation}
Janmitra project team. \emph{Prototype implementation and commit history}.
Repository revision \texttt{\Revision}, with major updates in \texttt{def3630}.
Application, tests, scripts, migrations, and data under \path{code/backend/};
CI configuration in \path{.github/workflows/backend.yml}.

\bibitem{architecture}
Janmitra project team. \emph{Architecture}.
\path{docs/architecture.md}, revision \texttt{\Revision}.

\bibitem{voice-design}
Janmitra project team. \emph{Browser Voice Architecture}.
\path{docs/voice-architecture.md}, revision \texttt{\Revision}.

\bibitem{runbook}
Janmitra project team. \emph{Run and Test Janmitra}.
\path{docs/run-and-test.md}, including the documented Gemini 3.1 provider-status
observations, revision \texttt{\Revision}.

\bibitem{catalogue-review}
Janmitra project team. \emph{Catalogue Review Checklist}.
\path{code/backend/data/REVIEW.md} and accompanying scheme JSON records,
revision \texttt{\Revision}.

\bibitem{source-check}
Janmitra repository. \emph{Automated Source Check: 9 October 2026}.
\path{code/backend/data/SOURCE-CHECK-2026-10-09.md}.
Automated research notes supplied for subsequent human review.

\bibitem{evaluation-fixtures}
Janmitra project team. \emph{Voice Evaluation Scenarios and Rubric}.
\path{code/backend/scripts/fixtures/voice-evaluation.json}, with
\path{code/backend/scripts/evaluate_voice_answer.py} and
\path{code/backend/scripts/smoke_voice.py}, revision \texttt{\Revision}.

\end{thebibliography}
\end{document}
